%
% Documento: Revisão Bibliografica
%
\chapter{Revisão da Literatura}\label{chap:revbibliografica}

\section{Faltas de Alta Impedância}

As faltas de alta impedância (FAIs) constituem um dos problemas mais persistentes e
difíceis enfrentados pela proteção de sistemas elétricos de distribuição
\cite{2, 3}. Caracterizam-se pelo contato não intencional entre um condutor
energizado e um meio de elevada resistência elétrica como solo seco, asfalto,
concreto, grama ou vegetação resultando em correntes de falta de baixa magnitude,
tipicamente na faixa de 0 a 85~A, insuficientes para acionar as proteções convencionais
por sobrecorrente, como fusíveis e relés de sobrecorrente \cite{1, Aljohani2020}.

Ao contrário das faltas de baixa impedância, que envolvem correntes elevadas e são
detectadas com facilidade pelas proteções tradicionais, as FAIs não representam risco
imediato de danos ao equipamento elétrico. Contudo, o condutor permanece energizado,
configurando sério risco à segurança pública e potencial de ignição de incêndios
\cite{3, Sedighizadeh2010}. A corrente de falta, próxima em magnitude à corrente normal
de carga, torna os relés de sobrecorrente praticamente cegos a esse tipo de evento
\cite{Aljohani2020}.

Do ponto de vista elétrico, as FAIs apresentam três características marcantes que
dificultam sua detecção e modelagem. A primeira é a não linearidade: o contato entre o
condutor e o meio resistivo gera um arco elétrico intermitente, cujo comportamento não
é descrito por uma resistência constante, mas por uma relação tensão--corrente não
linear e variável no tempo \cite{Elkalashy2007, WangCui2022}. A segunda é a assimetria
da forma de onda de corrente: o arco conduz de forma diferente nos semiciclos positivo e
negativo do ciclo da tensão, produzindo uma assimetria característica que é explorada
por vários métodos de detecção \cite{4, 5}. A terceira é o fenômeno de
\textit{buildup}: após o início da falta, a corrente cresce progressivamente ao longo
dos primeiros ciclos, refletindo o aquecimento e a evolução do ponto de contato
\cite{4}. Essas três características, combinadas, produzem assinaturas ricas em
componentes harmônicas e inter-harmônicas que se tornaram a base dos principais métodos
de detecção propostos na literatura \cite{Macedo2015, 7}.

A superfície de contato exerce influência direta sobre as propriedades elétricas da
falta. Materiais como asfalto seco e concreto apresentam resistência muito elevada,
produzindo correntes de falta próximas de zero; grama e solo úmido, por outro lado,
resultam em correntes ligeiramente maiores, porém ainda abaixo do limiar de atuação
convencional \cite{Aljohani2020, Sedighizadeh2010}. Em razão dessa dependência, as
assinaturas elétricas das FAIs variam consideravelmente conforme as condições ambientais
e o tipo de material em contato com o condutor \cite{Aljohani2020}.

Quando a superfície de contato é a vegetação, as FAIs assumem uma dimensão de risco
ainda mais elevada. O contato entre um condutor energizado e galhos ou copas de árvores
próximas cria um caminho de falta cujas propriedades elétricas dependem da espécie
vegetal, do teor de umidade e do estágio de desenvolvimento do contato
\cite{Ozansoy2020, 9}. Mesmo correntes de curtíssima duração e baixíssima magnitude
são capazes de carbonizar o material vegetal e originar brasas incandescentes,
configurando risco direto de ignição \cite{8, Ozansoy2020}. Estudos baseados em
ensaios controlados mostram que o processo de ignição evolui em estágios sequenciais, são eles: contato inicial, expulsão de umidade e carbonização. Vale ressaltar que cada um possui assinaturas
elétricas distintas na corrente de falta \cite{Ozansoy2020, Ozansoy2023}. A elevada
variabilidade entre espécies e condições ambientais torna a definição de limiares de
detecção confiáveis um problema em aberto, o que motivou a criação de programas
experimentais dedicados, como o \textit{Vegetation Conduction Ignition Test Program},
desenvolvido pelo governo de Victoria, Austrália \cite{vic_bushfire_report, 8}.

\section{Métodos de Detecção de FAIs}\label{sec:metodos}

A dificuldade de detecção das FAIs motivou, ao longo de décadas, o
desenvolvimento de uma ampla variedade de métodos. Como a corrente de falta não é
suficiente para atuar as proteções convencionais, as abordagens propostas buscam
explorar assinaturas elétricas indiretamente associadas ao fenômeno, como distorções
harmônicas, componentes de alta frequência e padrões não lineares da forma de onda
\cite{Sedighizadeh2010, Aljohani2020}. Apesar das décadas de pesquisa, ainda não existe
uma solução padronizada e universalmente aceita para o problema \cite{Sedighizadeh2010}.

Os trabalhos pioneiros remontam ao início da década de 1980. \citeonline{1} foram os
primeiros a propor sistematicamente o uso de componentes de alta frequência, na faixa
de 2 a 10~kHz, da corrente de alimentador como assinatura de FAI, construindo um
relé protótipo baseado em microcomputador testado em campo com resultados promissores.
Alguns anos depois, \citeonline{Huang1988} compararam quatro algoritmos baseados em
harmônicas de 2ª, 3ª e 5ª ordens em dados de faltas reais encenadas, estabelecendo
uma referência metodológica para avaliação comparativa de desempenho. \citeonline{3}
discutiram, além da eficácia técnica, as questões operacionais associadas à
implantação de detectores de FAI em serviço, reconhecendo que a aceitação pela
indústria dependia também de fatores não técnicos. Nesse período, Kim e Russell
\cite{KimRussell1989} introduziram o uso de raciocínio indutivo e sistemas especialistas
para classificação de FAIs, antecipando a transição para abordagens baseadas em
aprendizado.

A partir dos anos 2000, métodos baseados em processamento de sinais passaram a
dominar a literatura. \citeonline{5} propuseram o uso de árvores de decisão treinadas
com features harmônicas extraídas por simulação em EMTP, demonstrando que a
combinação de múltiplas harmônicas aumenta a confiabilidade da detecção.
\citeonline{Macedo2015} propuseram o uso de inter-harmônicos, gerados pela
variação do comprimento do arco elétrico, como assinatura de detecção, com
resultados validados em testes de campo em diferentes tipos de solo.
\citeonline{6} desenvolveram um algoritmo baseado na Transformada Wavelet Discreta
(DWT) capaz de identificar FAIs e indicar a área mais provável de ocorrência sem
exigir sincronismo de dados ou conhecimento dos parâmetros da rede. Já \citeonline{7}
utilizaram a Transformada de Fourier de Tempo Curto (STFT) para extrair magnitude e
fase das principais harmônicas da corrente, sendo capaz de distinguir FAIs de outros
distúrbios como chaveamento de capacitores e energização de alimentadores.

No campo das FAIs em vegetação, \citeonline{8} demonstraram que o conteúdo de alta
frequência dos sinais de tensão, associados ao arco elétrico intermitente típico
desse tipo de falta, constitui uma assinatura discriminativa eficaz, validada com
dados reais do programa experimental australiano. O classificador utilizado (\textit{boosted
decision trees}) alcançou alta sensibilidade tanto para faltas fase-terra quanto
fase-fase de baixa magnitude. A revisão mais recente sobre FAIs em árvores, conduzida
por \citeonline{9}, organiza as contribuições da área em três eixos:
modelagem, detecção e avaliação de risco de ignição e destaca a escassez de
dados reais como o principal gargalo para o avanço da área. Abordagens baseadas em
inteligência artificial, como SVM, redes neurais e XGBoost, têm sido progressivamente
aplicadas ao problema, sendo discutidas com maior profundidade na
Seção~\ref{sec:ia}.

\section{Técnicas de Processamento de Sinais}\label{sec:processamento}

As faltas de alta impedância produzem sinais de corrente de baixa amplitude,
não estacionários e fortemente contaminados pelo ruído de rede e pelas
variações normais de carga. Essas características tornam insuficiente a
análise direta no domínio do tempo ou o uso isolado de limiares de sobrecorrente,
exigindo técnicas de processamento de sinais capazes de revelar padrões
latentes nas formas de onda \cite{Sedighizadeh2010, Aljohani2020}. A escolha da
técnica influencia diretamente quais atributos podem ser extraídos e, portanto,
qual tipo de assinatura elétrica será utilizada na detecção ou
classificação da falta.

A Transformada de Fourier Discreta (DFT) tornou-se uma das primeiras ferramentas amplamente
empregadas na análise de FAIs, permitindo decompor o sinal de corrente em suas
componentes de frequência. Para um sinal discreto $x[n]$ com $N$ amostras, a
DFT é definida como

$$X[k] = \sum_{n=0}^{N-1} x[n]\, e^{-j2\pi kn/N}, \quad k = 0, 1, \ldots, N-1,$$

onde $X[k]$ representa a amplitude e fase da componente de frequência
$f_k = k f_s / N$, sendo $f_s$ a frequência de amostragem \cite{1, Huang1988}.
Na prática, a DFT é computada de forma eficiente pelo algoritmo FFT
(\textit{Fast Fourier Transform}). Contudo, por ser uma transformada global, a
DFT não captura variações temporais no conteúdo espectral, uma
limitação crítica para sinais transientes como os das FAIs.

A Transformada de Fourier de Tempo Curto (STFT) supera parcialmente essa
restrição ao multiplicar o sinal por uma função janela $w[n]$
deslizante antes de aplicar a DFT:

$$
\mathrm{STFT}\{x\}(m,k)
=
\sum_{n=-\infty}^{\infty}
x[n]\,w[n-m]\,
e^{-j\frac{2\pi}{N}kn}
$$

obtendo assim uma representação tempo-frequência de resolução
fixa. A escolha do tamanho da janela implica um compromisso: janelas longas
oferecem boa resolução em frequência mas baixa resolução temporal,
e vice-versa. \citeonline{7} demonstraram que a STFT permite extrair magnitude
e fase das harmônicas de baixa ordem com robustez suficiente para distinguir
FAIs de distúrbios comuns da rede, como chaveamento de bancos de capacitores e
energização de alimentadores, inclusive quando aplicada a oscilografias
reais.

A Transformada Wavelet Discreta (DWT) supera a limitação de resolução
fixa da STFT ao decompor o sinal em sub-bandas de frequência através de
filtros passa-alta $h[n]$ e passa-baixa $g[n]$ aplicados sucessivamente:

$$d_j[k] = \sum_{n} x[n]\, h[2k - n], \qquad a_j[k] = \sum_{n} x[n]\, g[2k - n],$$

onde $d_j$ são os coeficientes de detalhe (alta frequência) e $a_j$ os
coeficientes de aproximação (baixa frequência) no nível $j$ de
decomposição. Essa estrutura multirresolução permite analisar
transientes de curta duração com alta resolução temporal e, ao
mesmo tempo, capturar variações lentas com boa resolução em
frequência \cite{Elkalashy2007}. \citeonline{6} desenvolveram um algoritmo
baseado em DWT que monitora componentes de alta e baixa frequência de tensão
em múltiplos pontos do sistema, indicando a área mais provável de
ocorrência da falta sem exigir sincronismo de dados ou conhecimento dos
parâmetros do alimentador. No contexto das FAIs em vegetação,
\citeonline{Elkalashy2007} mostraram que re-ignições do arco elétrico
produzem transientes capturados com eficácia pelos coeficientes wavelet,
reforçando a utilidade da DWT como ferramenta de extração de atributos
nesse cenário.

Uma abordagem complementar e de grande relevância para a detecção de FAIs é o
uso de \textbf{filtros rejeita-faixa harmônica} para isolar as componentes
inter-harmônicas do sinal de corrente. A ideia central é que, durante uma FAI,
a variação estocástica do comprimento do arco elétrico gera componentes de
frequência que não são múltiplos inteiros da fundamental --- as chamadas
inter-harmônicas --- enquanto, em condições normais, o espectro de corrente é
dominado por harmônicas de ordem inteira e ruído de banda larga
\cite{Macedo2015}. Para revelar essas componentes, aplica-se ao sinal um banco
de filtros que rejeita as frequências harmônicas conhecidas
($f_0$, $2f_0$, $3f_0$, $\ldots$), preservando o resíduo inter-harmônico sobre
o qual são calculadas as métricas de detecção.

\citeonline{Macedo2015} incorporaram esse banco de filtros a relés de
religadores e seccionadores automáticos convencionais e validaram a metodologia
em testes de campo em diferentes tipos de solo, com medições instantâneas da
forma de onda de corrente durante a ocorrência de faltas reais. Os resultados
demonstraram que as inter-harmônicas constituem uma assinatura robusta do arco
intermitente das FAIs e que sua extração por filtros analíticos é compatível
com a infraestrutura já existente de proteção automática. Essa característica
torna o método particularmente atrativo do ponto de vista de implantação
industrial, pois não exige hardware dedicado de alta frequência nem
modificações estruturais no sistema de proteção \cite{Macedo2015}.

\section{Algoritmos de Inteligência Artificial}\label{sec:ia}

Os métodos baseados em processamento de sinais, embora eficazes na extração
de assinaturas espectrais, dependem frequentemente da definição manual de
limiares de decisão, cuja calibração é sensível às condições operacionais e ao
tipo de superfície de contato \cite{Aljohani2020, Sedighizadeh2010}. A
variabilidade inerente das FAIs, especialmente aquelas envolvendo
vegetação, cujas assinaturas dependem da espécie, umidade e estágio de
ignição, torna atrativa a substituição de regras fixas por modelos capazes
de aprender padrões diretamente dos dados. Esse contexto motivou a aplicação
crescente de algoritmos de inteligência artificial (IA) ao problema de
detecção e classificação de FAIs \cite{Aljohani2020}.

Uma abordagem intermediária nessa evolução é representada pelos métodos com
\textbf{limiar auto-adaptativo} (\textit{self-adaptive threshold}). Ao contrário
dos detectores com limiar fixo, esses métodos ajustam dinamicamente o critério de
decisão em função de propriedades estatísticas do sinal monitorado --- como média e
desvio padrão calculados em janelas deslizantes --- ou de condições operacionais da
rede, como variações de carga e topologia do alimentador. A premissa é que, durante
a operação normal, os indicadores extraídos do sinal apresentam comportamento
estacionário que serve de referência; a detecção ocorre quando uma mudança
estatisticamente significativa é observada em relação a essa linha de base
\cite{YuKhan1994, Sedighizadeh2010}.

Uma implementação concreta dessa ideia aplicada especificamente a FAIs é
apresentada por \cite{KimRussell1993}, que propõem um sistema de
detecção baseado em um modelo de elemento adaptativo. O sistema combina
múltiplos algoritmos de detecção operando em diferentes parâmetros de
frequência e integra seus resultados por meio de pesos associados a cada par
algoritmo-parâmetro, ajustados dinamicamente por uma fórmula elíptica de
correção. Diferentemente de um limiar fixo global, o critério de decisão
final emerge da combinação ponderada dos indicadores individuais, adaptando-se
ao comportamento do sinal ao longo do tempo. Validado em exemplos de
classificação falta/não-falta em alimentadores de distribuição, o método
demonstrou bom desempenho especialmente em segurança contra identificações
falsas \cite{KimRussell1993}.

Entre os métodos supervisionados, as árvores de decisão foram pioneiras nessa
transição. \citeonline{5} demonstraram que uma árvore treinada com
\textit{features} harmônicas simuladas em EMTP era capaz de distinguir FAIs de
chaveamentos e energizações com alta confiabilidade, ao custo de uma
frequência de amostragem modesta (1920~Hz). O uso de \textit{boosted decision
trees} que é uma combinação de múltiplas árvores com atualização ponderada dos erros, foi adotado por \citeonline{8} para classificação de VeHIFs a partir de
conteúdo de alta frequência, obtendo alta sensibilidade em dados reais
ruidosos. Métodos como SVM (\textit{Support Vector Machine}) também foram
aplicados ao problema, explorando a separabilidade das \textit{features} no
espaço transformado por núcleos não lineares \cite{Aljohani2020}.

Algoritmos de ensemble mais avançados, como Random Forest e XGBoost,
passaram a ser utilizados à medida que bases de dados reais tornaram-se
disponíveis. \citeonline{Ma2020} aplicaram um modelo denominado
\textit{Hybrid Step XGBoost} a 188 registros de faltas reais do programa
australiano PBSP, coletados em testes de campo com vegetação de copa
superior. O modelo atingiu 98,17\% de acurácia na classificação entre
contatos seguros e perigosos entre vegetação e linhas aéreas, superando redes
neurais profundas nos experimentos conduzidos. Esse resultado evidencia o
potencial de métodos baseados em gradiente aplicados diretamente a dados de
campo, sem recorrer a modelos simulados.

A disponibilidade pública do conjunto de dados VeHIF
\cite{Gomes2021VeHIF}, originado do programa experimental do governo de
Victoria, constitui um marco para a área, pois viabiliza a comparação objetiva
de algoritmos sob as mesmas condições operacionais reais
\cite{8, vic_bushfire_report}. A ausência histórica de bases públicas com
faltas reais em vegetação era apontada como um dos principais gargalos para o
avanço da área \cite{Gomes2021VeHIF}. O presente trabalho se insere nesse
contexto, propondo a extração e análise sistemática de atributos a partir
desse conjunto de dados com vistas à construção de um classificador
supervisionado.